% ====================================================================
% === SUPERSEDED (reader-facing) by Docs/notes/04_local_volatility.tex
%     (Vol-Fitter Technical Notes series). For the maintained,
%     comprehensive treatment with implementation + performance
%     appendices and code-generated figures, read that note.
% === RETAINED here as the equation-label citation source for code
%     docstrings, which reference this file by path and by its eq.
%     labels (e.g. "eq. q_logit", "note section 5.2").
% ====================================================================

\documentclass[11pt]{article}

\usepackage[a4paper,margin=1in]{geometry}
\usepackage{amsmath,amssymb,amsthm,mathtools}
\usepackage{bm}
\usepackage{booktabs,longtable,array}
\usepackage{graphicx}
\usepackage{enumitem}
\usepackage{hyperref}
\usepackage{xcolor}
\usepackage{microtype}
\usepackage{caption}
\usepackage{fancyhdr}

\hypersetup{
  colorlinks=true,
  linkcolor=blue!50!black,
  citecolor=blue!50!black,
  urlcolor=blue!50!black
}

\setlength{\headheight}{14pt}
\pagestyle{fancy}
\fancyhf{}
\lhead{Piecewise-affine local variance calibration}
\rhead{\thepage}

\newtheorem{definition}{Definition}
\newtheorem{proposition}{Proposition}
\newtheorem{remark}{Remark}
\newtheorem{lemma}{Lemma}

\newcommand{\R}{\mathbb{R}}
\newcommand{\E}{\mathbb{E}}
\newcommand{\1}{\mathbf{1}}
\newcommand{\dd}{\,\mathrm{d}}
\newcommand{\half}{\frac12}
\newcommand{\vloc}{\nu}
\newcommand{\Tmax}{T_{\max}}
\newcommand{\xmin}{x_{\min}}
\newcommand{\xmax}{x_{\max}}
\newcommand{\argmin}{\operatorname*{arg\,min}}

\title{\Large Calibration of Continuous Piecewise-Affine Local-Variance Surfaces\\[2mm]
\large from Sparse Vanilla Option and Variance-Swap Quotes}
\author{Technical note}
\date{\today}

\begin{document}
\maketitle

\begin{abstract}
This note describes a concrete calibration problem for local volatility: from a finite set
of option quotes across strikes and expiries, and possibly variance-swap quotes, choose a
continuous piecewise-affine local-variance surface on a triangulated time-moneyness grid.
The chosen surface is the one whose prices, obtained by the forward Dupire partial
differential equation, are closest to the quotes in a weighted least-squares sense. The
construction is inspired by the Andreasen-Huge philosophy of fitting local variance
directly through an implicit finite-difference pricing map, but the parameterization used
here is a continuous $P_1$ finite-element surface on an arbitrary triangulation.
The note derives the forward equation, the variance-swap identities, the discrete pricing
operator, sensitivities, an adjoint gradient, and a complete numerical example.
\end{abstract}

\tableofcontents

\section{Purpose and high-level idea}

A market volatility surface is observed only at finitely many points:
\[
  \{(T_j,K_j,\Pi_j^{\mathrm{mkt}}):j=1,\ldots,M_{\mathrm{opt}}\},
\]
where $T_j$ is expiry, $K_j$ is strike, and $\Pi_j^{\mathrm{mkt}}$ is a call or put
premium. There may also be variance-swap fair variance quotes
\[
  \{(S_q,V_q^{\mathrm{mkt}}):q=1,\ldots,M_{\mathrm{var}}\}.
\]
A smooth two-dimensional implied-volatility interpolation can look harmless while
introducing calendar-spread or butterfly arbitrage. Instead of interpolating implied
volatility, we fit a local variance function $\vloc(t,x)$ and price every quote through
a no-arbitrage diffusion model.

The method is:

\begin{enumerate}[label=\textbf{Step \arabic*.},leftmargin=2.2cm]
\item Convert all vanilla quotes to normalized forward call prices at normalized strike
      $x=K/F(0,T)$.
\item Choose a triangulation of a rectangle in $(t,x)$ whose number of vertices is of the
      same order as the number of quotes.
\item Parameterize the local variance by a continuous piecewise-affine function
      \[
        \vloc_\theta(t,x)=\sum_{\ell=1}^{m}\theta_\ell\phi_\ell(t,x),
      \]
      where $\phi_\ell$ are the usual hat functions of the triangulation.
\item For each parameter vector $\theta$, solve the forward Dupire PDE
      \[
        \partial_T c(T,x)=\frac12\,\vloc_\theta(T,x)x^2\partial_{xx}c(T,x),
        \qquad c(0,x)=(1-x)^+.
      \]
\item Minimize a weighted least-squares misfit to option quotes and variance-swap quotes,
      usually with positivity bounds and a mild roughness penalty.
\end{enumerate}

The direct calibration problem is nonlinear because the PDE solution depends nonlinearly
on $\theta$. It is nevertheless low-dimensional if $m$ is comparable to the quote count,
and each pricing evaluation is cheap because the one-dimensional Dupire equation leads
to tridiagonal linear systems.

\section{Market normalization}

\subsection{Deterministic rates and forwards}

Let $S_t$ be the spot price under a risk-neutral measure with deterministic short rate
$r(t)$ and deterministic dividend or foreign rate $q(t)$:
\[
  \frac{\dd S_t}{S_t}=(r(t)-q(t))\dd t+\sigma(t,S_t)\dd W_t.
\]
Let
\[
  D(0,T)=\exp\left(-\int_0^T r(s)\dd s\right),
  \qquad
  F(0,T)=S_0\exp\left(\int_0^T (r(s)-q(s))\dd s\right).
\]
Define the forward-normalized process
\[
  X_t=\frac{S_t}{F(0,t)}.
\]
Then $X_0=1$ and
\[
  \dd X_t=X_t\sqrt{\vloc(t,X_t)}\,\dd W_t,
  \qquad
  \vloc(t,x)=\sigma^2(t,F(0,t)x).
\]
Thus, after forward normalization, the model is a driftless martingale diffusion on
$(0,\infty)$:
\begin{equation}
  \boxed{\dd X_t=X_t\sqrt{\vloc(t,X_t)}\,\dd W_t,\qquad X_0=1.}
  \label{eq:forward_sde}
\end{equation}

\subsection{Normalized call prices}

Let $C^{\mathrm{mkt}}(T,K)$ and $P^{\mathrm{mkt}}(T,K)$ be domestic-currency option
prices at time $0$. The normalized forward call price is
\[
  c(T,x)=\frac{C(T,K)}{D(0,T)F(0,T)},\qquad x=\frac{K}{F(0,T)}.
\]
A put quote is converted to a call quote by put-call parity:
\[
  \frac{P(T,K)}{D(0,T)F(0,T)}=c(T,x)+x-1,
  \qquad
  c(T,x)=\frac{P(T,K)}{D(0,T)F(0,T)}+1-x.
\]
In the normalized variables, all vanilla call prices have the simple representation
\begin{equation}
  c(T,x)=\E[(X_T-x)^+].
  \label{eq:normalized_call_def}
\end{equation}
The boundary values are
\[
  c(T,0)=\E[X_T]=1,\qquad \lim_{x\to\infty}c(T,x)=0,
\]
and the initial condition is
\[
  c(0,x)=(1-x)^+.
\]

\section{Forward Dupire equation}

\subsection{Derivation from the forward density}

Assume first that $X_T$ has a smooth density $p(T,\cdot)$ and that integrations by parts
are justified. The density of the diffusion \eqref{eq:forward_sde} satisfies the forward
Kolmogorov equation
\begin{equation}
  \partial_T p(T,y)
  =\frac12\,\partial_{yy}\!\left(\vloc(T,y)y^2p(T,y)\right),
  \qquad p(0,y)=\delta_1(y).
  \label{eq:density_pde}
\end{equation}
From \eqref{eq:normalized_call_def},
\[
  c(T,x)=\int_x^\infty (y-x)p(T,y)\dd y.
\]
Differentiating with respect to strike gives
\[
  \partial_x c(T,x)=-\int_x^\infty p(T,y)\dd y,
  \qquad
  \partial_{xx}c(T,x)=p(T,x).
\]
Differentiating with respect to maturity and using \eqref{eq:density_pde},
\begin{align*}
  \partial_T c(T,x)
  &=\int_x^\infty (y-x)\partial_Tp(T,y)\dd y\\
  &=\frac12\int_x^\infty (y-x)
     \partial_{yy}\!\left(\vloc(T,y)y^2p(T,y)\right)\dd y\\
  &=\frac12\,\vloc(T,x)x^2p(T,x).
\end{align*}
Therefore
\begin{equation}
  \boxed{
  \partial_T c(T,x)=\frac12\,\vloc(T,x)x^2\partial_{xx}c(T,x),
  \qquad c(0,x)=(1-x)^+.
  }
  \label{eq:forward_dupire_normalized}
\end{equation}

This is the forward Dupire equation in normalized strike. The familiar local-variance
formula follows formally:
\begin{equation}
  \vloc(T,x)=\frac{2\,\partial_Tc(T,x)}{x^2\partial_{xx}c(T,x)}
  \qquad
  \text{where } \partial_{xx}c(T,x)>0.
  \label{eq:dupire_formula_normalized}
\end{equation}
The present calibration does not estimate the derivatives in
\eqref{eq:dupire_formula_normalized} from sparse data. Instead, it parameterizes
$\vloc$ and prices through \eqref{eq:forward_dupire_normalized}.

\subsection{Unnormalized form}

For reference, if $C(T,K)$ is the domestic discounted call price and rates/dividends are
deterministic, then
\begin{equation}
  \partial_T C
  =-q(T)C-(r(T)-q(T))K\partial_K C
    +\frac12\,\sigma^2(T,K)K^2\partial_{KK}C.
  \label{eq:dupire_unnormalized}
\end{equation}
The normalized equation \eqref{eq:forward_dupire_normalized} is usually preferable
numerically because it removes deterministic drift and discounting.

\subsection{No-arbitrage interpretation}

For a fixed maturity, call prices must be decreasing and convex in strike:
\[
  \partial_x c(T,x)\in[-1,0],
  \qquad
  \partial_{xx}c(T,x)\ge 0.
\]
Across maturities, normalized forward call prices should be nondecreasing:
\[
  \partial_T c(T,x)\ge 0.
\]
Equation \eqref{eq:forward_dupire_normalized} makes the last condition transparent:
if $\vloc(T,x)\ge0$ and $\partial_{xx}c(T,x)\ge0$, then $\partial_Tc(T,x)\ge0$.
At the continuous level, a positive local variance and the martingale diffusion
construction produce arbitrage-free European prices. A finite-difference implementation
should still monitor monotonicity and convexity, especially near boundaries and in wings.

\section{Piecewise-affine local variance on a triangulation}

\subsection{Triangulated parameter domain}

Fix a computational domain
\[
  \mathcal D=[0,\Tmax]\times[0,\xmax].
\]
A triangulation $\mathcal T$ is a finite collection of non-overlapping triangles whose
union is $\mathcal D$. Let the vertices be
\[
  z_\ell=(\tau_\ell,\xi_\ell),\qquad \ell=1,\ldots,m.
\]
The number $m$ is chosen to be of the same order as the number of available instruments.
For instance, if there are $18$ instruments, a grid with $20$ to $30$ vertices is a
reasonable first design.

Let $\phi_\ell$ be the nodal hat function associated with vertex $z_\ell$:
\[
  \phi_\ell(z_j)=\1_{\ell=j},\qquad
  \phi_\ell\text{ is affine on every triangle.}
\]
The $P_1$ finite-element space is
\[
  V_h=\{v:\mathcal D\to\R:\ v \text{ is continuous and affine on every triangle of }
  \mathcal T\}.
\]
The local variance surface is parameterized as
\begin{equation}
  \vloc_\theta(t,x)=\sum_{\ell=1}^{m}\theta_\ell\phi_\ell(t,x),
  \qquad \theta=(\theta_1,\ldots,\theta_m)^\top.
  \label{eq:p1_lv}
\end{equation}
If
\[
  0<\underline v\le \theta_\ell\le \overline v<\infty
  \qquad \text{for all }\ell,
\]
then $\underline v\le \vloc_\theta(t,x)\le\overline v$ everywhere on $\mathcal D$,
because a point inside a triangle is represented by nonnegative barycentric coordinates.

\subsection{Barycentric representation}

Let $z=(t,x)$ lie in a triangle with vertices $z_a,z_b,z_c$. Its barycentric coordinates
$(\lambda_a,\lambda_b,\lambda_c)$ are defined by
\[
  z=\lambda_a z_a+\lambda_b z_b+\lambda_c z_c,\qquad
  \lambda_a+\lambda_b+\lambda_c=1.
\]
Inside the triangle,
\[
  \lambda_a,\lambda_b,\lambda_c\ge0,
  \qquad
  \vloc_\theta(z)
  =\lambda_a\theta_a+\lambda_b\theta_b+\lambda_c\theta_c.
\]
Thus evaluation is cheap once each PDE grid point has been assigned to a triangle.

\subsection{Design of the triangulation}

A practical triangulation should include:

\begin{itemize}[leftmargin=1.1cm]
\item time vertices at or near important maturities;
\item moneyness vertices near liquid strikes, especially around the money;
\item left and right wing vertices for extrapolation;
\item a boundary $x=\xmax$ far enough that calls struck at $\xmax$ are negligible.
\end{itemize}

The surface is local variance, not implied variance. It need not visually resemble the
implied-volatility surface. Sparse data identify the local variance best in regions that
are reached with nonnegligible probability before quoted expiries. Boundary vertices and
far wings are therefore often controlled more by priors and regularization than by the
quotes themselves.

\section{Pricing map and finite-difference discretization}

\subsection{Spatial grid}

Let
\[
  0=x_0<x_1<\cdots<x_N=\xmax
\]
be a strike grid. Write $h_i=x_{i+1}-x_i$ and let $U_i^n$ approximate
$c(t_n,x_i)$, where $0=t_0<t_1<\cdots<t_{N_t}=\Tmax$.

For an interior point $x_i$, $1\le i\le N-1$, use the nonuniform central difference
\begin{equation}
  (D_{xx}U)_i
  =
  \frac{2}{h_{i-1}+h_i}
  \left(
    \frac{U_{i+1}-U_i}{h_i}
    -
    \frac{U_i-U_{i-1}}{h_{i-1}}
  \right).
  \label{eq:nonuniform_second_derivative}
\end{equation}
Define
\[
  a_i^-=\frac{x_i^2}{(h_{i-1}+h_i)h_{i-1}},
  \qquad
  a_i^+=\frac{x_i^2}{(h_{i-1}+h_i)h_i},
  \qquad
  a_i^0=-(a_i^-+a_i^+).
\]
Then
\[
  \frac12 x_i^2(D_{xx}U)_i
  =
  a_i^-U_{i-1}+a_i^0U_i+a_i^+U_{i+1}.
\]
With local variance value $\vloc_i^n=\vloc_\theta(t_n,x_i)$,
\[
  (A^n(\theta)U)_i
  =
  \vloc_i^n\left(a_i^-U_{i-1}+a_i^0U_i+a_i^+U_{i+1}\right).
\]
The matrix $A^n$ is tridiagonal on the interior unknowns.

\subsection{Implicit Euler step}

The boundary conditions are
\[
  U_0^n=1,\qquad U_N^n=0.
\]
The fully implicit step is
\begin{equation}
  \left(I-\Delta t_n A^{n+1}_{II}(\theta)\right)U_I^{n+1}
  =
  U_I^n+\Delta t_n b^{n+1}(\theta),
  \label{eq:implicit_step}
\end{equation}
where $I$ denotes interior nodes and $b^{n+1}$ contains the known boundary contribution.
For example, the first interior row receives
\[
  b^{n+1}_1=\vloc_\theta(t_{n+1},x_1)a_1^-\,U_0^{n+1}
          =\vloc_\theta(t_{n+1},x_1)a_1^-.
\]
All other boundary terms are zero if $U_N^{n+1}=0$.

The matrix
\[
  M_n(\theta)=I-\Delta t_n A^{n+1}_{II}(\theta)
\]
has positive diagonal entries, nonpositive off-diagonal entries, and is strictly
diagonally dominant under the standard boundary treatment. It is therefore an
$M$-matrix, so $M_n^{-1}\ge0$ componentwise. This is a useful stability and comparison
property of the implicit scheme.

\subsection{Higher-order time stepping: the generic two-level step}

The implicit step \eqref{eq:implicit_step} is first order in time, and the payoff kink
makes that order expensive. At the money the normalized call starts as
$c(t,1)\approx\sqrt{\vloc\,t/(2\pi)}$, so $\partial_tc\propto t^{-1/2}$ and a step of
length $\Delta t$ taken at time $t$ resolves the solution only to a relative accuracy of
order $\Delta t/t$. On a front expiry marched with $N$ uniform steps the resulting
implied-volatility error at the money is about $0.15\,\sigma/N$ whatever the expiry
(measured: $82$, $86$ and $217$ bp at $N=2$ on a 2-day SPY, a 27-day SPY and a 27-day
NVDA front; $6$, $8$ and $15$ bp at $N=32$), and the calibration then bends $\theta$ to
cancel it. A second-order scheme removes the $1/N$ term; the question is which one keeps
the stability of the implicit step.

Every time scheme the march supports is one instance of a single two-level relation on
the interior unknowns. With per-step coefficients
$(\gamma_n,\alpha_n,\beta_n,\varepsilon_n)$,
\begin{equation}
  \left(I-\gamma_n\Delta t_n A^{n+1}_{II}(\theta)\right)U_I^{n+1}
  =
  \alpha_n U_I^n-\beta_n U_I^{n-1}
  +\varepsilon_n\Delta t_n A^{n}_{II}(\theta)U_I^n
  +\Delta t_n\left(\gamma_n b^{n+1}(\theta)+\varepsilon_n b^{n}(\theta)\right),
  \label{eq:generic_two_level_step}
\end{equation}
where $b^{n}$ is the boundary vector of \eqref{eq:implicit_step} evaluated at level $n$.
The schemes are the rows of the coefficient table:

\begin{center}
\begin{tabular}{lcccc}
\toprule
scheme & $\gamma_n$ & $\alpha_n$ & $\beta_n$ & $\varepsilon_n$\\
\midrule
implicit Euler, eq.~\eqref{eq:implicit_step} & $1$ & $1$ & $0$ & $0$\\[1mm]
Crank--Nicolson (after the Rannacher start-up) & $\tfrac12$ & $1$ & $0$ & $\tfrac12$\\[1mm]
BDF2, eq.~\eqref{eq:bdf2_step} &
  $\dfrac{1+\omega_n}{1+2\omega_n}$ &
  $\dfrac{(1+\omega_n)^2}{1+2\omega_n}$ &
  $\dfrac{\omega_n^2}{1+2\omega_n}$ & $0$\\[3mm]
\bottomrule
\end{tabular}
\end{center}

The BDF2 row is the variable-step backward differentiation formula. Let
$\omega_n=\Delta t_n/\Delta t_{n-1}$ and take the quadratic through the three levels
$(t_{n-1},U^{n-1})$, $(t_n,U^n)$, $(t_{n+1},U^{n+1})$; its derivative at $t_{n+1}$ is
$\bigl[\tfrac{1+2\omega_n}{1+\omega_n}U^{n+1}-(1+\omega_n)U^n
+\tfrac{\omega_n^2}{1+\omega_n}U^{n-1}\bigr]/\Delta t_n$.
Setting it equal to the operator at the new level, $A^{n+1}U^{n+1}$, and multiplying by
$(1+\omega_n)/(1+2\omega_n)$ gives
\begin{equation}
  \gamma_n=\frac{1+\omega_n}{1+2\omega_n},\qquad
  \alpha_n=\frac{(1+\omega_n)^2}{1+2\omega_n},\qquad
  \beta_n=\frac{\omega_n^2}{1+2\omega_n},\qquad
  \varepsilon_n=0,
  \qquad
  \omega_n=\frac{\Delta t_n}{\Delta t_{n-1}}.
  \label{eq:bdf2_step}
\end{equation}
At $\omega_n=1$ these are the familiar $(\tfrac23,\tfrac43,\tfrac13,0)$.

Three facts follow from the table.
\begin{itemize}[leftmargin=1.1cm]
\item \emph{Constants are preserved.} $\alpha_n-\beta_n=1$ in every row (for BDF2,
      $(1+\omega_n)^2-\omega_n^2=1+2\omega_n$), and $A$ annihilates constants at interior
      rows since $a_i^0=-(a_i^-+a_i^+)$; a constant vector with matching boundary data is
      therefore a fixed point of every scheme, which is the consistency condition at
      order zero.
\item \emph{Order and damping.} The quadratic interpolant is exact for quadratics, so
      BDF2 has local error $O(\Delta t^3)$ and is second order, like Crank--Nicolson.
      They differ on the stiff modes. On $u'=\lambda u$ with $z=\lambda\Delta t$, the
      implicit step amplifies by $1/(1-z)$, Crank--Nicolson by $(1+z/2)/(1-z/2)$, and
      BDF2 by the roots of $(1-\gamma_nz)r^2-\alpha_nr+\beta_n=0$. As $z\to-\infty$ the
      implicit and BDF2 factors tend to $0$ (both are L-stable) whereas the
      Crank--Nicolson factor tends to $-1$. The stiffest modes of $A$ live on the finest
      strike cells, $|\lambda|\sim\vloc x^2/h^2$, and the payoff kink at $x=1$ loads them
      fully at $t=0$: Crank--Nicolson does not damp them but flips their sign at every
      step, which is the oscillation that makes it non-monotone on fine strike lattices.
      The Rannacher start-up (the first two steps implicit) damps the kink's high modes
      before Crank--Nicolson takes over; BDF2 needs no such repair.
\item \emph{The tridiagonal solve is unchanged.} The matrix on the left of
      \eqref{eq:generic_two_level_step} is $I-\gamma_n\Delta t_nA^{n+1}_{II}$, which for
      any $\gamma_n\in(0,1]$ has the sign pattern and strict diagonal dominance of $M_n$:
      it is an $M$-matrix with a nonnegative inverse, and the no-pivot tridiagonal solve
      of the implicit step applies to every scheme. What the two-level right-hand side
      does not carry over is the implicit step's comparison principle (the coefficient of
      $U^{n-1}$ is $-\beta_n<0$), so positivity of the marched density is checked on the
      output rather than inherited from the step; on every fitted surface of the
      measurement record of the subsection ``Graded grids'' the BDF2 reprice produced
      none negative.
\end{itemize}

BDF2 reads $U^{n-1}$, so the first step of a march is an implicit Euler step. Growth of
the step is limited: at $z=0$ the characteristic polynomial factors as
$(r-1)(r-\beta_n)$, and the parasitic root $\beta_n=\omega_n^2/(1+2\omega_n)$ stays below
$1$ exactly when $\omega_n<1+\sqrt2$, the zero-stability bound of variable-step BDF2.
The march therefore restarts with an implicit Euler step whenever a step grows by more
than a factor $2$ (a margin under that bound). Shrinking steps never restart:
$\omega_n\to0$ sends $(\gamma_n,\alpha_n,\beta_n)\to(1,1,0)$ and the step tends to
implicit Euler smoothly. On the graded time grid of the subsection ``Graded grids'' the
growth never exceeds $1+c=1.25$, so a BDF2 march restarts only at its first step.

\subsection{Observation operator}

Let $R_j$ be the row vector that linearly interpolates grid values at strike $x_j$.
If $T_j$ coincides with a time step $t_{n(j)}$, then the model quote is
\[
  P_j(\theta)=R_jU^{n(j)}(\theta).
\]
If it does not, one may either include $T_j$ in the time grid or interpolate between
adjacent time levels. In calibration, it is usually best to force all quoted expiries to
be time steps.

\subsection{Graded grids}

\paragraph{Time grid.}
The estimate $\partial_tc\propto t^{-1/2}$ at the money says that the solution's own
time scale is $t$: a step $\Delta t$ taken at time $t$ resolves the kink to a relative
accuracy of order $\Delta t/t$, so equal accuracy per step means geometric growth,
$\Delta t_n\le c\,t_n$ with $c=0.25$, from a first step $\Delta t_0$ equal to $1\%$ of the
first mark. Geometric growth alone is not enough. The local variance $\vloc_\theta$ is
piecewise affine in $t$ with kinks at the vertex rows $t_i$, so $\partial_tc$ has a kink
and $\partial_{tt}c$ a jump at every row; a step straddling a row loses the scheme's
order, and the slabs between rows are exactly where a purely geometric grid has grown
coarse (measured: $12$--$22$ bp at intermediate expiries). Every listed expiry
\emph{and} every vertex row is therefore a grid point (a \emph{mark}), each slab between
consecutive marks receives at least $8$ steps, and an absolute ceiling $\Delta t\le0.05$
closes the rule:
\[
  \Delta t_n=\min\!\left(c\,t_n,\ 0.05,\ (1+c)\,\Delta t_{n-1},\ \tfrac18\,\text{slab}\right),
\]
a mark being hit exactly by splitting the slab's remainder evenly once fewer than three
steps remain. A step consequently never grows by more than the factor $1+c$, so the BDF2
march of \eqref{eq:bdf2_step} restarts only at its first step. The shape of the grid does
not depend on the scheme; the implicit legacy keeps its per-interval uniform rule
($\Delta t\le0.01$, any interval that would receive fewer than $8$ steps lifted to $32$)
so that its results are unchanged bit for bit.

\paragraph{Strike lattice.}
A uniform lattice must take the shortest expiry's step everywhere: a 2-day front needs a
cell of about $1/800$ near the money, so a uniform lattice carries some $1700$ nodes over
$[0,\xmax]$ for every expiry, though that front's density lives within a few percent of
$x=1$. The graded lattice grants each expiry its own region and step. With
$\sigma_j\sqrt{\tau_j}$ the at-the-money standard deviation of expiry $j$, its region
$[x_{\mathrm{lo}}^j,x_{\mathrm{hi}}^j]$ is its traded strike range widened to
$\pm6\,\sigma_j\sqrt{\tau_j}$ in log-moneyness (beyond which the density is below
$e^{-18}$ of its peak), and its step is
\[
  h_j=\operatorname{clip}\!\left(0.15\,\sigma_j\sqrt{\tau_j},\ \tfrac1{800},\ 0.01\right),
\]
the resolution its own quotes require. Outside every region the step is the wing step
$h_{\mathrm{wing}}=0.02$. The union of these requirements is a step function, and a
lattice cannot follow it abruptly: the truncation error of the nonuniform stencil
\eqref{eq:nonuniform_second_derivative} is
$\tfrac13(h_i-h_{i-1})\,\partial_{xxx}U+O(h^2)$, first order in general, second order
in practice only while consecutive cells differ by a small factor. The requirements are
therefore smoothed into the Lipschitz envelope
\[
  h(x)=\min\Bigl(h_{\mathrm{wing}},\ \min_j\bigl[h_j+(r-1)\operatorname{dist}(x,\text{region}_j)\bigr]\Bigr),
  \qquad r=1.15,
\]
the largest step function below every requirement whose slope is at most $r-1$; walking
with step $h(x)$ then changes the cell by at most a factor of about $r$ between
neighbours, and the leading error term above is bounded by $(r-1)h/3$. The lattice is
integrated outward from $x=1$ in both directions, which makes $x=1$ a node by
construction --- the static replication \eqref{eq:variance_swap_static_replication}
splits its put and call legs there and the at-the-money row of the triangulation sits
there, both of which need it exactly --- and closes on $0$ and $\xmax$ with a last cell
between one half and one and a half steps.

\paragraph{Measured effect.}
On the fitted surfaces of the measurement record (a Bloomberg SPY and NVDA chain, a
true-weekly capture and a SPY daily ladder), BDF2 on the graded time grid reads at most
$5$ bp of operator error per expiry, repriced against a reference marched with $256$
steps per interval, at $98$--$116$ steps, where the uniform per-interval implicit rule
marched $51$--$271$ steps and left $15$--$170$ bp per expiry for the calibration to
absorb. With a second-order time scheme the strike lattice is the residual --- its own
floor at the front is $5$--$9$ bp rms at the $0.15\,\sigma\sqrt\tau$ step, $2.5$--$4$
with the step halved --- which is what the graded lattice keeps fine where the density
lives and coarse where nothing does: a SPY surface carrying a 2-day rung drops from
about $1700$ to about $400$ strike nodes at the same near-the-money resolution.

\section{Variance swaps}

\subsection{Expected realized variance}

For the diffusion \eqref{eq:forward_sde}, the continuously sampled total variance over
$[0,T]$ is
\[
  I(T)=\E\left[\int_0^T \vloc(t,X_t)\dd t\right].
\]
The annualized fair variance rate is
\[
  K_{\mathrm{var}}(T)=\frac{I(T)}{T}.
\]
A variance-swap market quote should therefore be compared either as a rate
$K_{\mathrm{var}}$ or as total variance $I(T)=TK_{\mathrm{var}}(T)$. In least squares,
total variance is often numerically convenient because it scales additively in time.

\subsection{Log-contract identity}

Ito's formula gives
\[
  \dd\log X_t
  =
  -\frac12\vloc(t,X_t)\dd t+\sqrt{\vloc(t,X_t)}\dd W_t.
\]
Taking expectations,
\begin{equation}
  I(T)
  =
  \E\left[\int_0^T\vloc(t,X_t)\dd t\right]
  =
  -2\,\E[\log X_T].
  \label{eq:log_contract_identity}
\end{equation}
For any twice differentiable payoff $g$, the static replication formula around anchor
$\kappa>0$ is
\[
  \E[g(X_T)]
  =
  g(\kappa)+g'(\kappa)(\E[X_T]-\kappa)
  +\int_0^\kappa g''(k)P(T,k)\dd k
  +\int_\kappa^\infty g''(k)C(T,k)\dd k,
\]
where $C(T,k)=\E[(X_T-k)^+]$ and $P(T,k)=\E[(k-X_T)^+]$.
Taking $\kappa=1$, $\E[X_T]=1$, and $g(x)=-2\log x$ yields
\begin{equation}
  \boxed{
  I(T)=2\int_0^1\frac{P(T,k)}{k^2}\dd k
      +2\int_1^\infty\frac{C(T,k)}{k^2}\dd k.
  }
  \label{eq:variance_swap_static_replication}
\end{equation}
In a numerical implementation the integrals are truncated and the wings are extrapolated.
If the truncation is material, the variance-swap quote is partly a statement about wing
assumptions, not only about observed listed options.

\subsection{Source PDE alternative}

Instead of using \eqref{eq:variance_swap_static_replication}, define
\[
  g(t,x)=\E_{t,x}\left[\int_t^T\vloc(s,X_s)\dd s\right].
\]
Then $g$ satisfies the backward source PDE
\begin{equation}
  \partial_tg(t,x)+\frac12\vloc(t,x)x^2\partial_{xx}g(t,x)+\vloc(t,x)=0,
  \qquad
  g(T,x)=0.
  \label{eq:variance_swap_source_pde}
\end{equation}
The model total variance is $I(T)=g(0,1)$. This PDE is useful when variance swaps are
included but one does not want their value to depend on a numerical quadrature of the
computed option surface.

\section{Calibration objective}

\subsection{Weighted least squares}

Let $y_j$ be the normalized market call price of option quote $j$, and let
$\eta_j>0$ be its price tolerance, typically derived from bid-ask width. Let
$z_q$ be a total variance quote $S_qK_{\mathrm{var}}^{\mathrm{mkt}}(S_q)$ with tolerance
$\zeta_q>0$. Let $P_j(\theta)$ be the model option price and $Z_q(\theta)$ the model
total variance.

The basic calibration problem is
\begin{equation}
  \boxed{
  \theta^\star
  =
  \argmin_{\underline v\le\theta\le\overline v}
  \left[
  \frac12\sum_{j=1}^{M_{\mathrm{opt}}}
    \left(\frac{P_j(\theta)-y_j}{\eta_j}\right)^2
  +
  \frac12\sum_{q=1}^{M_{\mathrm{var}}}
    \left(\frac{Z_q(\theta)-z_q}{\zeta_q}\right)^2
  +
  \frac{\lambda}{2}\|L(\theta-\theta^{\mathrm{ref}})\|_2^2
  \right].
  }
  \label{eq:calibration_objective}
\end{equation}
Here $L$ is optional. It may be a graph-gradient, a second-difference matrix, or a
finite-element roughness matrix. The vector $\theta^{\mathrm{ref}}$ is a prior, often a
flat variance or the previous day's calibrated surface.

The phrase ``closest to the quotes'' is precisely defined by the choice of tolerances
$\eta_j,\zeta_q$ and regularization. With bid and ask prices $b_j,a_j$, a common choice is
\[
  y_j=\frac{a_j+b_j}{2},
  \qquad
  \eta_j=\frac{a_j-b_j}{2}
\]
after normalization. If the bid-ask width is unrealistically small, a floor should be
added.

\subsection{Existence and uniqueness}

If the feasible set
\[
  \Theta=\{\theta\in\R^m:\underline v\le\theta_\ell\le\overline v\}
\]
is compact and the finite-difference pricing map is continuous in $\theta$, then
\eqref{eq:calibration_objective} has at least one minimizer. Uniqueness is not guaranteed.
The inverse problem is usually underdetermined in wings and between maturities, and the
objective is nonlinear. Regularization should therefore be viewed as part of the model
selection rule, not merely as a numerical convenience.

\subsection{Exact fit versus robust fit}

Even if $m$ is close to the number of quotes, an exact fit may fail for several reasons:
the quotes may contain arbitrage, bid-ask mid prices may not be mutually consistent, the
triangulation may be too coarse, or boundary extrapolation may matter. In production, it
is preferable to fit within bid-ask and keep a stable local-variance surface rather than
to force exact interpolation of noisy mids.

\section{Sensitivities and gradients}

\subsection{Continuous sensitivity PDE}

Let
\[
  s_\ell(T,x)=\frac{\partial c(T,x;\theta)}{\partial\theta_\ell}.
\]
Since
\[
  \vloc_\theta(t,x)=\sum_{\ell=1}^m\theta_\ell\phi_\ell(t,x),
\]
differentiating \eqref{eq:forward_dupire_normalized} gives
\begin{equation}
  \partial_T s_\ell
  =
  \frac12\vloc_\theta(T,x)x^2\partial_{xx}s_\ell
  +
  \frac12\phi_\ell(T,x)x^2\partial_{xx}c,
  \qquad
  s_\ell(0,x)=0.
  \label{eq:sensitivity_pde}
\end{equation}
Thus each parameter sensitivity is the solution of the same parabolic operator with a
source term proportional to the option gamma.

For variance swaps computed through the source PDE \eqref{eq:variance_swap_source_pde},
the sensitivity $h_\ell=\partial g/\partial\theta_\ell$ satisfies
\begin{equation}
  \partial_th_\ell
  +\frac12\vloc_\theta x^2\partial_{xx}h_\ell
  +\frac12\phi_\ell x^2\partial_{xx}g
  +\phi_\ell=0,
  \qquad h_\ell(T,x)=0.
  \label{eq:var_sensitivity_pde}
\end{equation}

\subsection{Discrete forward sensitivities}

Differentiate the implicit step
\[
  M_n(\theta)U^{n+1}=U^n+\Delta t_n b^{n+1}(\theta).
\]
Ignoring the notationally minor boundary derivative, the sensitivity matrix
$S^n=\partial U^n/\partial\theta\in\R^{(N-1)\times m}$ satisfies
\begin{equation}
  M_n(\theta)S^{n+1}
  =
  S^n
  +
  \Delta t_n
  \left(\frac{\partial A^{n+1}}{\partial\theta}U^{n+1}\right),
  \qquad
  S^0=0.
  \label{eq:discrete_sensitivity}
\end{equation}
More explicitly, for parameter $\theta_\ell$ and row $i$,
\[
  \left[
  \left(\frac{\partial A^{n+1}}{\partial\theta_\ell}\right)U^{n+1}
  \right]_i
  =
  \phi_\ell(t_{n+1},x_i)
  \left(a_i^-U_{i-1}^{n+1}+a_i^0U_i^{n+1}+a_i^+U_{i+1}^{n+1}\right).
\]
Since the same tridiagonal matrix $M_n$ is used for all $m$ sensitivities, all columns of
$S^{n+1}$ can be solved simultaneously by a tridiagonal solve with multiple right-hand
sides.

The generic step \eqref{eq:generic_two_level_step} differentiates the same way. The
operator is linear in $\theta$ and the coefficients do not depend on it, so
\begin{equation}
  \left(I-\gamma_n\Delta t_nA^{n+1}\right)S^{n+1}
  =
  \alpha_nS^n-\beta_nS^{n-1}
  +\varepsilon_n\Delta t_nA^nS^n
  +\gamma_n\Delta t_n\left(\frac{\partial A^{n+1}}{\partial\theta}U^{n+1}\right)
  +\varepsilon_n\Delta t_n\left(\frac{\partial A^{n}}{\partial\theta}U^{n}\right),
  \qquad S^0=S^{-1}=0,
  \label{eq:generic_sensitivity_step}
\end{equation}
with the row-wise source of the previous display at each level and the boundary
derivative again omitted; for implicit Euler this is \eqref{eq:discrete_sensitivity}.
The cost per level is read off the terms. The tridiagonal solve is shared by all $m$
columns as before. With $\varepsilon_n=0$ the source keeps the implicit step's single
term, evaluated at the new level, and the only addition is the combination
$\alpha_nS^n-\beta_nS^{n-1}$, one vector operation per column. BDF2's sensitivity march
therefore costs an implicit one plus that combination, whereas Crank--Nicolson
($\varepsilon_n=\tfrac12$) adds the stencil $A^nS^n$ on every column and a second source
at the old level, roughly doubling the work per level. This is why BDF2, not
Crank--Nicolson, is the default second-order scheme.

\subsection{Discrete adjoint gradient}

When $m$ is large, an adjoint is more efficient. Suppose the objective contributes a
gradient $g^n=\partial\Phi/\partial U^n$ at each time step, including interpolation rows
for quotes observed at $t_n$. Initialize adjoints $\bar U^n=g^n$. For
$n=N_t-1,\ldots,0$:

\[
  \text{solve}\qquad M_n(\theta)^\top z^{n+1}=\bar U^{n+1},
\]
then update
\begin{align}
  \nabla_\theta\Phi
  &\mathrel{+}=
  \Delta t_n
  \left[
  z^{n+1\top}
  \left(\frac{\partial A^{n+1}}{\partial\theta_1}U^{n+1}\right),
  \ldots,
  z^{n+1\top}
  \left(\frac{\partial A^{n+1}}{\partial\theta_m}U^{n+1}\right)
  \right]^\top,                                      \label{eq:adjoint_grad}\\
  \bar U^n&\mathrel{+}=z^{n+1}.
\end{align}
Boundary derivative terms are added in the same manner. The adjoint requires one forward
pricing run and one backward sweep, independent of $m$.

\section{Algorithm}

A robust implementation is as follows.

\begin{enumerate}[label=\arabic*.,leftmargin=0.9cm]
\item \textbf{Normalize market data.}
      Convert all calls and puts to normalized forward calls. Convert variance-swap
      rates to total variance.
\item \textbf{Pre-clean quotes.}
      Remove or deweight obviously crossed, stale, or arbitrage-inconsistent quotes.
\item \textbf{Build triangulation.}
      Choose vertices in $(t,x)$ and precompute barycentric weights
      $\phi_\ell(t_n,x_i)$ at all PDE grid points.
\item \textbf{Choose bounds and prior.}
      For example, $\underline v=(1\%)^2$ and $\overline v=(200\%)^2$ in extreme equity
      markets, or tighter product-specific values.
\item \textbf{Price for a trial $\theta$.}
      Solve \eqref{eq:implicit_step}, or the generic step \eqref{eq:generic_two_level_step}
      with the BDF2 coefficients \eqref{eq:bdf2_step} on the graded grids; record $U^n$
      at quote maturities.
\item \textbf{Compute residuals.}
      Apply interpolation rows for option quotes. Compute variance-swap model values by
      either \eqref{eq:variance_swap_static_replication} or
      \eqref{eq:variance_swap_source_pde}.
\item \textbf{Compute derivatives.}
      Use forward sensitivities if $m$ is small; use the adjoint if $m$ is large.
\item \textbf{Optimize.}
      Use a bound-constrained least-squares method or a smooth transformation such as
      \[
        \theta_\ell
        =
        \underline v+
        \frac{\overline v-\underline v}{1+\exp(-\alpha_\ell)}.
      \]
\item \textbf{Validate.}
      Check price residuals, implied-volatility residuals, call monotonicity,
      convexity in strike, calendar monotonicity, and stability under small data
      perturbations.
\end{enumerate}

\section{Detailed numerical example}

This example uses normalized forwards, so $F(0,T)=1$ and $D(0,T)=1$. The quote strikes
are therefore moneyness values.

\subsection{Synthetic quote set}

The synthetic local-variance surface is generated by sampling
\begin{equation}
  \vloc_\star(t,x)
  =
  0.032+0.006t+0.030(1-x)^2+0.012(1-x)
  +0.004\sin(\pi t)\exp\!\left[-\left(\frac{x-1}{0.35}\right)^2\right]
  \label{eq:example_true_surface}
\end{equation}
at the triangulation nodes and then using piecewise-affine interpolation. This gives a
mild left skew and increasing term level.

The finite set of option quotes is five strikes at three maturities. The following table
shows normalized call prices and Black-Scholes implied volatilities.

\begin{table}[h!]
\centering
\caption{Synthetic option quote set. Implied volatility is annualized and shown in percent.}
\begin{tabular}{rrrr}
\toprule
$T$ & $x$ & call price & implied vol. (\%)\\
\midrule
0.25 & 0.80 & 0.200277 & 19.071\\
0.25 & 0.90 & 0.105645 & 18.579\\
0.25 & 1.00 & 0.036544 & 18.327\\
0.25 & 1.10 & 0.007310 & 18.245\\
0.25 & 1.20 & 0.000861 & 18.280\\
0.50 & 0.80 & 0.202596 & 19.285\\
0.50 & 0.90 & 0.115765 & 19.017\\
0.50 & 1.00 & 0.053085 & 18.832\\
0.50 & 1.10 & 0.019104 & 18.701\\
0.50 & 1.20 & 0.005456 & 18.616\\
1.00 & 0.80 & 0.211163 & 19.630\\
1.00 & 0.90 & 0.133968 & 19.411\\
1.00 & 1.00 & 0.076657 & 19.245\\
1.00 & 1.10 & 0.039690 & 19.126\\
1.00 & 1.20 & 0.018833 & 19.061\\
\bottomrule
\end{tabular}
\label{tab:quote_set}
\end{table}

The variance-swap quotes are generated from the same model by the log-contract identity,
with numerical integration on $x\in[0.01,2.20]$.

\begin{table}[h!]
\centering
\caption{Synthetic variance-swap quotes. Rates are annualized variances, not volatilities.}
\begin{tabular}{rr}
\toprule
$T$ & market variance rate\\
\midrule
0.25 & 0.033931\\
0.50 & 0.035713\\
1.00 & 0.037479\\
\bottomrule
\end{tabular}
\label{tab:var_quotes}
\end{table}

\subsection{Triangulation and numerical grid}

The local-variance triangulation has
\[
  \tau\in\{0,0.5,1.0\},\qquad
  \xi\in\{0,0.70,0.90,1.00,1.10,1.30,2.20\}.
\]
Thus there are $m=21$ local-variance parameters. The number of instruments is
$15+3=18$, so the local-variance grid has the same order of size as the quote set.

The Dupire PDE is solved on
\[
  x_i=0.01i,\quad i=0,\ldots,220,
  \qquad
  t_n=0.005n,\quad n=0,\ldots,200.
\]
The calibration starts from the flat initial guess $\theta_\ell=0.04$ and uses bounds
\[
  0.005\le \theta_\ell\le0.20.
\]
The objective uses option price tolerance $2\times10^{-4}$, total variance tolerance
$2\times10^{-4}$, and a mild second-difference roughness penalty.

\subsection{Calibrated local-variance values}

The calibrated nodal local variances are:

\begin{table}[h!]
\centering
\caption{Calibrated nodal local variances $\theta_{\ell}$.}
\resizebox{\textwidth}{!}{%
\begin{tabular}{rrrrrrrr}
\toprule
$t\backslash x$ & 0.00 & 0.70 & 0.90 & 1.00 & 1.10 & 1.30 & 2.20\\
\midrule
0.00 & 0.04607 & 0.03973 & 0.03348 & 0.03198 & 0.03109 & 0.03429 & 0.04105\\
0.50 & 0.05064 & 0.04220 & 0.04006 & 0.03905 & 0.03769 & 0.03613 & 0.04515\\
1.00 & 0.05513 & 0.04675 & 0.03963 & 0.03813 & 0.03699 & 0.03755 & 0.04621\\
\bottomrule
\end{tabular}}
\label{tab:calibrated_variances}
\end{table}

Equivalently, the calibrated nodal local volatilities, in percent, are:

\begin{table}[h!]
\centering
\caption{Calibrated nodal local volatilities $100\sqrt{\theta_{\ell}}$.}
\resizebox{\textwidth}{!}{%
\begin{tabular}{rrrrrrrr}
\toprule
$t\backslash x$ & 0.00 & 0.70 & 0.90 & 1.00 & 1.10 & 1.30 & 2.20\\
\midrule
0.00 & 21.46 & 19.93 & 18.30 & 17.88 & 17.63 & 18.52 & 20.26\\
0.50 & 22.50 & 20.54 & 20.01 & 19.76 & 19.41 & 19.01 & 21.25\\
1.00 & 23.48 & 21.62 & 19.91 & 19.53 & 19.23 & 19.38 & 21.50\\
\bottomrule
\end{tabular}}
\label{tab:calibrated_vols}
\end{table}

The far-left and far-right nodes are not tightly identified by the option quotes. Their
values are influenced by the variance-swap quotes, the boundary conditions, and the
smoothness penalty. This is typical: the wings of a local-variance surface are model
selection choices unless liquid wing instruments are present.

\subsection{Fit to options}

The table below compares market and model normalized call prices. A price bp means
$10^{-4}$ in normalized forward price. The last column is the implied-volatility error in
basis points of volatility.

\begin{longtable}{rrrrrr}
\caption{Fit to option quotes.}\label{tab:fit_options}\\
\toprule
$T$ & $x$ & market call & model call & price error (bp) & vol error (bp)\\
\midrule
\endfirsthead
\toprule
$T$ & $x$ & market call & model call & price error (bp) & vol error (bp)\\
\midrule
\endhead
0.25 & 0.80 & 0.200277 & 0.200282 & +0.050 & +4.30\\
0.25 & 0.90 & 0.105645 & 0.105645 & -0.002 & -0.02\\
0.25 & 1.00 & 0.036544 & 0.036544 & -0.003 & -0.02\\
0.25 & 1.10 & 0.007310 & 0.007308 & -0.018 & -0.15\\
0.25 & 1.20 & 0.000861 & 0.000867 & +0.054 & +1.79\\
0.50 & 0.80 & 0.202596 & 0.202595 & -0.011 & -0.16\\
0.50 & 0.90 & 0.115765 & 0.115763 & -0.014 & -0.07\\
0.50 & 1.00 & 0.053085 & 0.053088 & +0.030 & +0.11\\
0.50 & 1.10 & 0.019104 & 0.019099 & -0.048 & -0.21\\
0.50 & 1.20 & 0.005456 & 0.005461 & +0.051 & +0.43\\
1.00 & 0.80 & 0.211163 & 0.211185 & +0.225 & +1.21\\
1.00 & 0.90 & 0.133968 & 0.133967 & -0.012 & -0.04\\
1.00 & 1.00 & 0.076657 & 0.076666 & +0.086 & +0.22\\
1.00 & 1.10 & 0.039690 & 0.039689 & -0.012 & -0.03\\
1.00 & 1.20 & 0.018833 & 0.018846 & +0.125 & +0.45\\
\bottomrule
\end{longtable}

The root-mean-square normalized price error is
\[
  7.56\times10^{-6},
\]
and the largest absolute normalized price error is
\[
  2.25\times10^{-5}.
\]
The root-mean-square implied-volatility error is $1.26$ volatility basis points.

\subsection{Fit to variance swaps}

\begin{table}[h!]
\centering
\caption{Fit to variance-swap quotes. Rate error is in variance basis points.}
\begin{tabular}{rrrr}
\toprule
$T$ & market rate & model rate & rate error (bp)\\
\midrule
0.25 & 0.033931 & 0.033943 & +0.116\\
0.50 & 0.035713 & 0.035710 & -0.029\\
1.00 & 0.037479 & 0.037453 & -0.258\\
\bottomrule
\end{tabular}
\label{tab:fit_var}
\end{table}

The variance swaps are matched to within roughly $0.3$ variance basis points in this
example. Tightening the variance-swap tolerance would improve those rows at the possible
expense of vanilla option fit or surface smoothness.

\subsection{Interpretation}

Several features are worth noting.

\begin{enumerate}[leftmargin=0.9cm]
\item The calibration does not interpolate implied volatility. The implied-volatility
      smile is an output after the local-variance PDE has been solved.
\item The local-variance surface is not uniquely determined in regions that the process
      rarely visits. Regularization and bounds are economically meaningful inputs.
\item Variance swaps are especially informative about the average variance level, because
      they constrain $E[\int_0^T\vloc(t,X_t)\dd t]$ rather than a payoff localized around
      one strike.
\item The forward Dupire PDE turns a global surface-fitting problem into repeated
      one-dimensional tridiagonal solves, so moderately sized calibrations are fast.
\end{enumerate}

\section{Implementation details and safeguards}

\subsection{Choice of residual units}

Normalizing all options by $D(0,T)F(0,T)$ keeps residuals dimensionless. A tolerance of
$10^{-4}$ then means one basis point of forward price. Deep out-of-the-money options may
require vega-weighted or implied-volatility residuals to avoid being dominated by tiny
premiums. A robust production objective often combines price residuals for liquid options
with floors in implied-volatility or vega units.

\subsection{Arbitrage checks}

Before calibration, check the input prices at each maturity:
\[
  c(T,x)\ge(1-x)^+,\qquad
  c(T,x)\le1,\qquad
  c \text{ decreasing and convex in }x.
\]
Across maturities, check
\[
  c(T_2,x)\ge c(T_1,x)\quad\text{for }T_2>T_1
\]
after interpolation to common strikes. If the raw quote set violates these inequalities,
a local-volatility model cannot fit every mid exactly.

After calibration, check the computed model prices on a dense grid, not only at quoted
strikes. A finite-difference approximation can introduce small numerical violations even
when the continuous model is arbitrage-free.

\subsection{Boundary and wing treatment}

The boundary $x=\xmax$ should be chosen so that calls at $\xmax$ are negligible at all
calibrated maturities. The lower boundary $x=0$ is natural for calls because
$c(T,0)=1$. Variance-swap replication is much more sensitive to wings because of the
$k^{-2}$ weight in \eqref{eq:variance_swap_static_replication}. If variance swaps are
used, wing extrapolation must be explicit and stable.

\subsection{Regularization choices}

Common choices of $L$ in \eqref{eq:calibration_objective} include:

\[
  \|L\theta\|^2
  =
  \sum_{\text{edges }(a,b)}w_{ab}(\theta_a-\theta_b)^2
\]
for graph-gradient smoothing, or second-difference penalties in time and moneyness:
\[
  \sum_{i,j}\left(\theta_{i,j-1}-2\theta_{i,j}+\theta_{i,j+1}\right)^2
  +
  \rho\sum_{i,j}\left(\theta_{i-1,j}-2\theta_{i,j}+\theta_{i+1,j}\right)^2.
\]
A prior penalty
\[
  \|\theta-\theta^{\mathrm{ref}}\|^2
\]
is useful for daily calibration stability. The prior may be yesterday's surface.

\subsection{Parameter bounds}

Bounds prevent pathological local variances. They also guarantee compactness of the
feasible set and therefore existence of a minimizer. In practice the lower bound should
be positive but small enough not to bias low-volatility products. The upper bound should
be high enough not to bind in normal markets but low enough to prevent optimizer
excursions in poorly identified wings.

\section{Extensions}

\subsection{Multiple underlyings or smiles sharing parameters}

For a set of underlyings with related smiles, one can fit several local-variance surfaces
with shared regularization or common factors. The pricing PDE remains one-dimensional for
each underlying.

\subsection{American or barrier instruments}

The calibrated local variance can be used in backward pricing PDEs or Monte Carlo for
path-dependent instruments. The finite set of European quotes fixes the marginal laws of
$X_T$ only to the extent that they are fitted; path-dependent prices remain model
dependent.

\subsection{Alternative loss functions}

Least squares is convenient but not mandatory. Robust losses, bid-ask constrained
calibration, or maximum-likelihood-style objectives can be substituted:
\[
  b_j\le P_j(\theta)\le a_j
\]
for as many quotes as possible, plus a roughness minimization, is often more meaningful
than exact matching of mids.

\section{Conclusion}

A continuous piecewise-affine local-variance surface on a triangulated grid gives a
flexible but controlled way to fit sparse option and variance-swap quotes. The forward
Dupire PDE provides the pricing map from local variance to European option prices, while
the log-contract identity or a source PDE incorporates variance swaps. Positivity of the
nodal variances gives positivity of the entire $P_1$ surface, and an implicit
finite-difference method (implicit Euler, or second-order BDF2 on grids graded from the
payoff kink and from each expiry's own width) gives stable tridiagonal solves. The final calibration is a
bound-constrained nonlinear least-squares problem whose regularization should be treated
as part of the model specification.

\appendix

\section{Proof of the static replication identity}

Let $g\in C^2((0,\infty))$ and choose an anchor $\kappa>0$. For any $x>0$,
Taylor's formula with integral remainder can be written as
\[
  g(x)
  =
  g(\kappa)+g'(\kappa)(x-\kappa)
  +\int_0^\kappa g''(k)(k-x)^+\dd k
  +\int_\kappa^\infty g''(k)(x-k)^+\dd k.
\]
Taking expectations and using $\E[X_T]=1$ gives
\[
  \E[g(X_T)]
  =
  g(\kappa)+g'(\kappa)(1-\kappa)
  +\int_0^\kappa g''(k)P(T,k)\dd k
  +\int_\kappa^\infty g''(k)C(T,k)\dd k.
\]
With $\kappa=1$ and $g(x)=-2\log x$, $g(1)=g'(1)(1-1)=0$ and
$g''(k)=2/k^2$, which gives \eqref{eq:variance_swap_static_replication}.

\section{Why nodal positivity implies surface positivity}

Let $z$ be in a triangle with vertices $z_a,z_b,z_c$. Then
\[
  z=\lambda_az_a+\lambda_bz_b+\lambda_cz_c,
  \qquad
  \lambda_a,\lambda_b,\lambda_c\ge0,\qquad
  \lambda_a+\lambda_b+\lambda_c=1.
\]
For a $P_1$ function,
\[
  \vloc_\theta(z)=\lambda_a\theta_a+\lambda_b\theta_b+\lambda_c\theta_c.
\]
If every nodal value satisfies $\underline v\le \theta_\ell\le\overline v$, then
\[
  \underline v
  =
  (\lambda_a+\lambda_b+\lambda_c)\underline v
  \le
  \vloc_\theta(z)
  \le
  (\lambda_a+\lambda_b+\lambda_c)\overline v
  =
  \overline v.
\]
Thus positivity and upper bounds need only be imposed at vertices.

\section{Minimal pseudocode}

\begin{verbatim}
Input:
    normalized option quotes (T_j, x_j, y_j, eta_j)
    total variance quotes (S_q, z_q, zeta_q)
    triangulation nodes z_l and triangles
    PDE grids {t_n}, {x_i}
    bounds v_min, v_max

Precompute:
    barycentric weights Phi[n,i,l] = phi_l(t_n, x_i)
    interpolation rows R_j for option strikes

Objective(theta):
    U = payoff (1 - x)^+
    for n = 0,...,Nt-1:
        v_i = sum_l theta_l Phi[n+1,i,l]
        build tridiagonal M_n = I - dt A(v)
        solve M_n U_new = U + boundary_terms
        U = U_new
        store U if t_{n+1} is a quote maturity

    residuals = []
    for each option quote j:
        residuals.append((R_j U(T_j) - y_j)/eta_j)

    for each variance quote q:
        compute Z_q(theta) from log-contract quadrature
        residuals.append((Z_q(theta) - z_q)/zeta_q)

    residuals.extend( sqrt(lambda) * L(theta - theta_ref) )
    return residuals

Solve:
    minimize 0.5 * ||Objective(theta)||^2
    subject to v_min <= theta_l <= v_max
\end{verbatim}

\begin{thebibliography}{9}

\bibitem{AndreasenHuge}
J.~Andreasen and B.~N. Huge.
\newblock Volatility interpolation.
\newblock \emph{Risk}, March 2011. Earlier version available as SSRN working
paper, 2010.

\bibitem{Dupire}
B.~Dupire.
\newblock Pricing with a smile.
\newblock \emph{Risk}, 7(1):18--20, 1994.

\bibitem{BlackScholes}
F.~Black and M.~Scholes.
\newblock The pricing of options and corporate liabilities.
\newblock \emph{Journal of Political Economy}, 81(3):637--654, 1973.

\bibitem{Gatheral}
J.~Gatheral.
\newblock \emph{The Volatility Surface: A Practitioner's Guide}.
\newblock Wiley, 2006.

\bibitem{ContTankov}
R.~Cont and P.~Tankov.
\newblock \emph{Financial Modelling with Jump Processes}.
\newblock Chapman and Hall/CRC, 2004.

\end{thebibliography}

\end{document}
